\begin{table}[htbp]
\centering
\small
\renewcommand{\arraystretch}{0.90}
\caption{Granger Directional Precedence: Master Nominal Core System (Chile, 1960--1980, $N=250$)}
\label{tab:granger_nominal_core}
\resizebox{\textwidth}{!}{%
\begin{tabular}{llccccc}
\toprule
\multicolumn{7}{l}{\textbf{Panel A: Primary Baseline at SBIC-Selected Lag ($k^*$)}} \\
\midrule
\textbf{Dependent Variable (DV)} & \textbf{Predictor (INDV)} & \textbf{Lag ($k^*$)} & \textbf{$F$-Statistic} & \textbf{$p$-value} & \textbf{$\sum \hat{\beta}_i$} & \textbf{Granger Conclusion} \\
\midrule
Base Money Growth ($g_{H,t}$) & Inflation ($\pi_t$) & 3 & 20.318$^{***}$ & $< 0.0001$ & +0.595 & $\pi \to g_H$ detected \\
Inflation ($\pi_t$) & Base Money Growth ($g_{H,t}$) & 3 & 5.544$^{***}$ & 0.0011 & +0.421 & $g_H \to \pi$ detected \\
\midrule
\multicolumn{7}{l}{\textbf{Panel B: Lag Sensitivity and Horizon Persistence ($k = 1, \dots, 4$)}} \\
\midrule
\textbf{Hypothesis / Direction} & \textbf{$k = 1$} & \textbf{$k = 2$} & \textbf{$k = 3$} & \textbf{$k = 4$} & \multicolumn{2}{c}{\textbf{Persistence Summary}} \\
\midrule
Inflation ($\pi_t$) $\to$ Base Money Growth ($g_{H,t}$) & 28.32$^{***}$ & 18.52$^{***}$ & 20.32$^{***}$$^a$ & 15.00$^{***}$ & \multicolumn{2}{c}{Robust at all horizons ($p < 0.05$)} \\
Base Money Growth ($g_{H,t}$) $\to$ Inflation ($\pi_t$) & 14.56$^{***}$ & 9.28$^{***}$ & 5.54$^{***}$$^a$ & 2.72$^{**}$ & \multicolumn{2}{c}{Robust at all horizons ($p < 0.05$)} \\
\bottomrule
\end{tabular}%
}
\begin{minipage}{0.96\textwidth}
\vspace{0.3em}
\footnotesize \textbf{Notes:} Monthly series spanning 1960:01--1980:12 ($N=250$). Evaluated via standard Granger causality tests in stationary VAR($k$) specifications. Panel~A shows that at the SBIC-selected lag ($k^*=3$), Granger causality is bidirectional: price inflation predicts base money growth ($\pi \to g_H$, $F = 20.318, p < 0.0001, \sum \hat{\beta}_i = +0.595$), while base money also Granger-causes inflation ($g_H \to \pi$, $F = 5.544, p = 0.0011, \sum \hat{\beta}_i = +0.421$). Panel~B shows that inflation Granger-causes base-money growth more strongly and more persistently across the tested lag structure ($k=1..4$), an empirical ordering consistent with defensive monetary accommodation. Panel~A reports the primary baseline at the SBIC-selected optimal lag ($k^*$, denoted by $^a$). Panel~B reports finite-sample $F$-statistics across tested horizons $k = 1, \dots, 4$. Significance: $^{***} p < 0.01$, $^{**} p < 0.05$, $^{*} p < 0.10$. In accordance with the paper's methodological contract, Granger causality denotes incremental predictive precedence within a stationary VAR and does not by itself establish structural causal identification.
\end{minipage}
\end{table}
